
\section{Appendix D. Negative Results}

We document failed experiments to help future work avoid dead ends.

\subsection{D.1 Attempts That Failed}

\subsubsection{D.1.1 Trainable SVD Base}

\textbf{Goal.} Let $U, V$ train alongside adapters.

\textbf{Result.} Base factors drifted during recovery, undoing the compression. 
Accuracy peaked early then collapsed.

\textbf{Conclusion.} Keep base frozen (\texttt{requires\\_grad=False}).

\subsubsection{D.1.2 Very Low Energy Threshold ($\tau$=0.80)}

\textbf{Goal.} Aggressive compression.

\textbf{Result.} Accuracy dropped to {{TAU80\_ACC}}\%.

\textbf{Conclusion.} $\tau$=0.80 too aggressive for attention weights.

\subsubsection{D.1.3 Very High Threshold ($\tau$=0.99)}

\textbf{Goal.} Near-lossless compression.

\textbf{Result.} Accuracy was near-baseline but compression ratio was only 
{{TAU99\_RATIO}}$\times$.

\textbf{Conclusion.} $\tau$=0.99 defeats the purpose.

\subsubsection{D.1.4 No Residual Adapter (r\_L=0)}

\textbf{Goal.} Pure truncation without adapter.

\textbf{Result.} Accuracy dropped to {{NO\_ADAPT\_ACC}}\%.

\textbf{Conclusion.} Adapters are essential.

\subsubsection{D.1.5 Both A and B Random}

\textbf{Goal.} Different initialization.

\textbf{Result.} Initial forward pass differed from pretrained; accuracy 
degraded before recovery.

\textbf{Conclusion.} Zero-init one of the factors.

\subsubsection{D.1.6 Both A and B Zero}

\textbf{Goal.} Symmetric initialization.

\textbf{Result.} Gradient for both zero; no learning.

\textbf{Conclusion.} One must be nonzero.

\subsection{D.2 Abandoned Variants}

\subsubsection{D.2.1 Per-Head Rank}

We tested separate rank per head. Result: no improvement, 4$\times$ slower 
(SVD per head).

\subsubsection{D.2.2 Learned $\tau$}

Trained per-layer $\tau$ during recovery. Result: weak gradient signal, no 
improvement.

\subsubsection{D.2.3 Per-Projection Threshold}

Different $\tau$ for Q, K, V, O. Result: no clear pattern; uniform $\tau$=0.92 
works.

\subsubsection{D.2.4 Diagonal-Energy Rank}

Thresholded $|\sigma_i|$ instead of $\sigma_i^2$. Result: different rank, 
slightly worse accuracy.

\subsubsection{D.2.5 Fixed Rank Higher}

Tested fixed rank 192. Result: better accuracy than adaptive at $\tau$=0.92, 
but worse compression. Trade-off.

\subsection{D.3 Failed Hyperparameters}

\subsubsection{Learning Rate}

\begin{table}[H]
\centering
\begin{tabular}{lll}
\toprule
\textbf{LR} & \textbf{DistilBERT} & \textbf{Loss stability} \\
\midrule
5e-5 & {{D\_LR5e5}} & Stable \\
1e-4 & {{D\_LR1e4}} & Stable \\
2e-4 & {{D\_LR2e4}} & Stable \\
5e-4 & {{D\_LR5e4}} & Unstable \\
1e-3 & NaN & NaN \\
\bottomrule
\end{tabular}
\end{table}

\subsubsection{Batch Size}

\begin{table}[H]
\centering
\begin{tabular}{llll}
\toprule
\textbf{BS} & \textbf{DistilBERT} & \textbf{RoBERTa} & \textbf{GPT-2} \\
\midrule
2 & {{D\_BS2}} & {{R\_BS2}} & {{G\_BS2}} \\
4 & {{D\_BS4}} & {{R\_BS4}} & {{G\_BS4}} \\
8 & {{D\_BS8}} & {{R\_BS8}} & {{G\_BS8}} \\
16 & {{D\_BS16}} & OOM & OOM \\
\bottomrule
\end{tabular}
\end{table}

\subsubsection{Residual Rank}

\begin{table}[H]
\centering
\begin{tabular}{lll}
\toprule
\textbf{$r_L$} & \textbf{DistilBERT} & \textbf{Note} \\
\midrule
1 & {{D\_RL1}} & Underfit \\
2 & {{D\_RL2}} & Underfit \\
4 & {{D\_RL4}} & Acceptable \\
8 & {{D\_RL8}} & Good \\
16 & {{D\_RL16}} & Best \\
32 & {{D\_RL32}} & Diminishing returns \\
64 & {{D\_RL64}} & Overkill \\
\bottomrule
\end{tabular}
\end{table}

\subsection{D.4 Fixes That Didn't Help}

\begin{table}[H]
\centering
\begin{tabular}{lll}
\toprule
\textbf{Fix} & \textbf{Problem} & \textbf{Result} \\
\midrule
Orthogonal adapter init & Slow convergence & No change \\
Gradient clipping & Stability & No change \\
Adaptive $\tau$ per layer & Rank quality & No improvement \\
Post-hoc base fine-tune & Accuracy & Drifted from compression \\
Adding dropout to adapters & Overfitting & No improvement \\
\bottomrule
\end{tabular}
\end{table}

\subsection{D.5 Architecture-Specific Failures}

\subsubsection{D.5.1 T5 Verification Incomplete}

T5 support was partially tested. Full verification on encoder and decoder 
blocks was deferred due to time constraints.

\subsubsection{D.5.2 Qwen Partial}

Qwen3.5 with GQA was partially tested. Full verification pending.

\subsection{D.6 Multi-Seed Analysis}

\begin{table}[H]
\centering
\begin{tabular}{llll}
\toprule
\textbf{Seed} & \textbf{DistilBERT} & \textbf{RoBERTa} & \textbf{GPT-2} \\
\midrule
42 & {{SEED42\_D}} & {{SEED42\_R}} & {{SEED42\_G}} \\
123 & {{SEED123\_D}} & {{SEED123\_R}} & {{SEED123\_G}} \\
456 & {{SEED456\_D}} & {{SEED456\_R}} & {{SEED456\_G}} \\
**Mean** & **{{SEED\_MEAN\_D}}** & **{{SEED\_MEAN\_R}}** & **{{SEED\_MEAN\_G}}** \\
**Std** & **{{SEED\_STD\_D}}** & **{{SEED\_STD\_R}}** & **{{SEED\_STD\_G}}** \\
\bottomrule
\end{tabular}
\end{table}

\textbf{Observation.} Standard deviation ~0.15\% across seeds. Single-seed results 
are representative.

\subsection{D.7 Lessons Learned}

1. \textbf{Freeze base factors} -- never train them.
2. \textbf{Zero-init one adapter factor} -- both zero or both random fails.
3. \textbf{$\tau$=0.92 is robust} -- no need to tune extensively.
4. \textbf{$r_L=16$ is sufficient} -- diminishing returns above.
5. \textbf{No calibration is a feature} -- enables zero-data settings.
6. \textbf{Per-layer rank matters} -- fixed rank loses ~1.5\% accuracy.

\subsection{D.8 Recommendations}

Based on our negative results:

1. Start with default hyperparameters ($\tau$=0.92, r\_L=16).
2. Tune $\tau$ only if accuracy is unacceptable.
3. Increase r\_L if recovery plateaus below target.
4. Keep base frozen.
5. Verify on one model before full sweep.

\subsection{D.9 Unresolved Questions}

1. \textbf{Optimal $\tau$ schedule:} Could per-layer $\tau$ beat uniform?
2. \textbf{Combined with A3/AASVD:} Would activation-aware rank beat 
   energy-thresholded?
3. \textbf{Interaction with quantization:} INT8 + Adaptive SVD?
4. \textbf{Learning rate schedule:} Would cosine scheduling help?
5. \textbf{Multi-task calibration:} Does rank transfer across tasks?